\begin{ex}
Xăng-ti-mét viết tắt là chữ gì trong bảng chữ cái tiếng Anh? Số chữ cái viết tắt là bao nhiêu?
\shortans{2}
\loigiai{
Xăng-ti-mét được viết tắt là cm (gồm 2 chữ cái c và m).
\textbf{Đáp số:} $2$.
}
\end{ex}

\begin{ex}
Tính: $8\text{ cm} + 2\text{ cm}$. Kết quả là bao nhiêu xăng-ti-mét?
\shortans{10}
\loigiai{
Ta có: $8\text{ cm} + 2\text{ cm} = 10\text{ cm}$.
\textbf{Đáp số:} $10$.
}
\end{ex}

\begin{ex}
Tính: $10\text{ cm} - 3\text{ cm}$. Kết quả là bao nhiêu xăng-ti-mét?
\shortans{7}
\loigiai{
Ta có: $10\text{ cm} - 3\text{ cm} = 7\text{ cm}$.
\textbf{Đáp số:} $7$.
}
\end{ex}

\begin{ex}
Tính: $12\text{ cm} + 5\text{ cm}$. Kết quả là bao nhiêu xăng-ti-mét?
\shortans{17}
\loigiai{
Ta có: $12\text{ cm} + 5\text{ cm} = 17\text{ cm}$.
\textbf{Đáp số:} $17$.
}
\end{ex}

\begin{ex}
Tính: $19\text{ cm} - 7\text{ cm}$. Kết quả là bao nhiêu xăng-ti-mét?
\shortans{12}
\loigiai{
Ta có: $19\text{ cm} - 7\text{ cm} = 12\text{ cm}$.
\textbf{Đáp số:} $12$.
}
\end{ex}

\begin{ex}
Điền số thích hợp: $5\text{ cm} + \dots\text{ cm} = 9\text{ cm}$.
\shortans{4}
\loigiai{
Số cần điền là: $9 - 5 = 4$.
\textbf{Đáp số:} $4$.
}
\end{ex}

\begin{ex}
Điền số thích hợp: $15\text{ cm} - \dots\text{ cm} = 10\text{ cm}$.
\shortans{5}
\loigiai{
Số cần điền là: $15 - 10 = 5$.
\textbf{Đáp số:} $5$.
}
\end{ex}

\begin{ex}
Một chiếc bút sáp màu dài 8 cm, chiếc bút lông dài 14 cm. Hỏi chiếc bút lông dài hơn chiếc bút sáp bao nhiêu xăng-ti-mét?
\shortans{6}
\loigiai{
Bút lông dài hơn bút sáp là: $14 - 8 = 6$ (cm).
\textbf{Đáp số:} $6$.
}
\end{ex}

\begin{ex}
Đoạn thẳng AB dài 6 cm, đoạn thẳng BC dài 4 cm. Ba điểm A, B, C thẳng hàng theo thứ tự đó. Hỏi đoạn thẳng AC dài bao nhiêu xăng-ti-mét?
\shortans{10}
\loigiai{
Độ dài đoạn thẳng AC là: $6 + 4 = 10$ (cm).
\textbf{Đáp số:} $10$.
}
\end{ex}

\begin{ex}
Một sợi dây đồng dài 18 cm. Bạn Hùng cắt lấy 8 cm để làm thí nghiệm. Hỏi sợi dây đồng còn lại dài bao nhiêu xăng-ti-mét?
\shortans{10}
\loigiai{
Độ dài sợi dây đồng còn lại là: $18 - 8 = 10$ (cm).
\textbf{Đáp số:} $10$.
}
\end{ex}

\begin{ex}
Tính: $7\text{ cm} + 3\text{ cm} + 4\text{ cm}$. Kết quả là bao nhiêu xăng-ti-mét?
\shortans{14}
\loigiai{
Thực hiện: $7 + 3 = 10$; $10 + 4 = 14$ (cm).
\textbf{Đáp số:} $14$.
}
\end{ex}

\begin{ex}
Tính: $18\text{ cm} - 5\text{ cm} - 3\text{ cm}$. Kết quả là bao nhiêu xăng-ti-mét?
\shortans{10}
\loigiai{
Thực hiện: $18 - 5 = 13$; $13 - 3 = 10$ (cm).
\textbf{Đáp số:} $10$.
}
\end{ex}

\begin{ex}
Một cây thước dài 20 cm bị gãy mất một đầu dài 4 cm. Hỏi đoạn thước còn lại dài bao nhiêu xăng-ti-mét?
\shortans{16}
\loigiai{
Đoạn thước còn lại dài: $20 - 4 = 16$ (cm).
\textbf{Đáp số:} $16$.
}
\end{ex}

\begin{ex}
Con sâu đo bò được 5 cm trên chiếc lá, rồi bò thêm 6 cm nữa đến cuống lá. Hỏi con sâu đo đã bò được tất cả bao nhiêu xăng-ti-mét?
\shortans{11}
\loigiai{
Tổng quãng đường sâu đo bò là: $5 + 6 = 11$ (cm).
\textbf{Đáp số:} $11$.
}
\end{ex}

\begin{ex}
Một hộp bút dài 17 cm, một cục tẩy dài 3 cm. Hỏi cần bao nhiêu xăng-ti-mét nữa để cục tẩy dài bằng hộp bút?
\shortans{14}
\loigiai{
Hiệu độ dài là: $17 - 3 = 14$ (cm).
\textbf{Đáp số:} $14$.
}
\end{ex}

\begin{ex}
Đoạn thẳng MN dài 13 cm, đoạn thẳng PQ ngắn hơn đoạn MN là 3 cm. Hỏi đoạn thẳng PQ dài bao nhiêu xăng-ti-mét?
\shortans{10}
\loigiai{
Độ dài đoạn thẳng PQ là: $13 - 3 = 10$ (cm).
\textbf{Đáp số:} $10$.
}
\end{ex}

\begin{ex}
Tổng độ dài của hai đoạn thẳng MN và PQ ở bài toán trên là bao nhiêu xăng-ti-mét?
\shortans{23}
\loigiai{
Tổng độ dài là: $13 + 10 = 23$ (cm).
\textbf{Đáp số:} $23$.
}
\end{ex}

\begin{ex}
Một băng giấy dài 15 cm. Bạn Mai cắt băng giấy thành hai đoạn, trong đó có một đoạn dài 8 cm. Hỏi đoạn còn lại dài bao nhiêu xăng-ti-mét?
\shortans{7}
\loigiai{
Đoạn băng giấy còn lại dài: $15 - 8 = 7$ (cm).
\textbf{Đáp số:} $7$.
}
\end{ex}

\begin{ex}
Tính: $20\text{ cm} - 10\text{ cm} + 5\text{ cm}$.
\shortans{15}
\loigiai{
Thực hiện: $20 - 10 = 10$; $10 + 5 = 15$ (cm).
\textbf{Đáp số:} $15$.
}
\end{ex}

\begin{ex}
Một con cào cào nhảy bước thứ nhất được 8 cm, bước thứ hai nhảy được 9 cm. Hỏi cả hai bước cào cào nhảy được bao nhiêu xăng-ti-mét?
\shortans{17}
\loigiai{
Cả hai bước cào cào nhảy được: $8 + 9 = 17$ (cm).
\textbf{Đáp số:} $17$.
}
\end{ex}

\begin{ex}
Chiều dài chiếc điện thoại di động là 14 cm, chiều rộng là 7 cm. Chiều dài hơn chiều rộng bao nhiêu xăng-ti-mét?
\shortans{7}
\loigiai{
Chiều dài hơn chiều rộng là: $14 - 7 = 7$ (cm).
\textbf{Đáp số:} $7$.
}
\end{ex}

\begin{ex}
Tìm số thích hợp: $11\text{ cm} + \dots\text{ cm} = 16\text{ cm}$.
\shortans{5}
\loigiai{
Số cần điền là: $16 - 11 = 5$.
\textbf{Đáp số:} $5$.
}
\end{ex}

\begin{ex}
Tìm số thích hợp: $19\text{ cm} - \dots\text{ cm} = 11\text{ cm}$.
\shortans{8}
\loigiai{
Số cần điền là: $19 - 11 = 8$.
\textbf{Đáp số:} $8$.
}
\end{ex}

\begin{ex}
Một cây nến cao 16 cm. Sau một giờ cháy bớt đi 6 cm. Hỏi cây nến còn lại cao bao nhiêu xăng-ti-mét?
\shortans{10}
\loigiai{
Chiều cao còn lại của cây nến là: $16 - 6 = 10$ (cm).
\textbf{Đáp số:} $10$.
}
\end{ex}

\begin{ex}
Đoạn dây thép thứ nhất dài 12 cm, đoạn dây thép thứ hai dài 7 cm. Nối hai đoạn lại (không tính mối nối), sợi dây mới dài bao nhiêu xăng-ti-mét?
\shortans{19}
\loigiai{
Chiều dài sợi dây mới là: $12 + 7 = 19$ (cm).
\textbf{Đáp số:} $19$.
}
\end{ex}
